\subsection{Rounded weights in the marginal-moment interval}

Fix exponents $0<\ell_{\rm s}<1$, $0<\kappa<\mu<1-\ell_{\rm s}$,
$\nu>0$, and $0<e<1$ with $\kappa\le e$. For integer $n\ge1$, use
the scales of \cite[Section~9, lines~32--41]{OpenAI2026AreaLaw}:
\begin{align}
    L & =\lfloor n^{1-\ell_{\rm s}}\rfloor, \\
    m & =\lfloor n^\mu\rfloor,              \\
    K & =\lfloor L/(8m)\rfloor,             \\
    D & =\lceil n^\kappa\rceil.
    \label{eq:al_moment_rounded_scales}
\end{align}
Set $\alpha=1-\ell_{\rm s}-\mu>0$. Only these three exponents enter
(\ref{eq:al_moment_rounded_scales}).

\begin{theorem}[Quantitative growth of the integer band count]
    \label{thm:al_moment_band_power}
    \lean{TNLean.PEPS.AreaLaw.Scan.ScannerExponents.exists_eventually_rpow_le_mul_K}
    \leanok
    There is a constant $C_K\ge1$ such that all sufficiently large
    integers $n$ satisfy $n\ge2$, $K\ge1$, and $n^\alpha\le C_KK$.
\end{theorem}
\begin{proof}\leanok
    \uses{lem:al_scanner_scale_facts}
    The numerical scale estimates give positive $m,K,D$ and
    $nD^2/(mK)\le C_Kn^{\ell_{\rm s}}D^2$.
    Cancel $D^2$, multiply by $mK$, and use $m\le n^\mu$ to obtain
    \begin{align}
        n & \le C_Kn^{\ell_{\rm s}}mK
        \le C_KK n^{\ell_{\rm s}+\mu}.
    \end{align}
    Dividing by the positive last power gives the assertion.
\end{proof}

\begin{theorem}[Square root of the rounded marginal budget]
    \label{thm:al_moment_sqrt_budget}
    \lean{TNLean.PEPS.AreaLaw.Scan.ScannerExponents.sqrt_budget_le}
    \leanok
    Let $n\ge1$, $C_B\ge0$ and $B\in\mathbb R$ satisfy
    $B\le C_BnD(\log n)^{12}$. Then
    \begin{align}
        \sqrt B & \le\sqrt{2C_B}\,n^{(1+\kappa)/2}(\log n)^6.
    \end{align}
    The real square root is zero on negative arguments.
\end{theorem}
\begin{proof}\leanok
    Since $n^\kappa\ge1$, rounding gives
    $D\le n^\kappa+1\le2n^\kappa$. Thus
    $B\le2C_Bn^{1+\kappa}(\log n)^{12}$. The right side is the
    square of the asserted nonnegative upper bound.
\end{proof}

\begin{theorem}[One threshold for every bounded weight and budget]
    \label{thm:al_moment_weight_smallness}
    \lean{TNLean.PEPS.AreaLaw.Scan.ScannerExponents.eventually_mul_a_sqrt_budget_le,
        TNLean.PEPS.AreaLaw.Scan.ScannerExponents.exists_a_le_moment_radius}
    \leanok
    Fix $\omega\in\mathbb R$, $C_B,A\ge0$ and $c>0$. Assume the strict
    margin
    \begin{align}
        \delta & =\alpha-\omega-(1+\kappa)/2>0.
        \label{eq:al_moment_margin}
    \end{align}
    There is an integer $N\ge2$, chosen before $n,W,B$, such that every
    $n\ge N$ has $K\ge1$ and the following property. For every
    $0\le W\le n^\omega$ and real $B\le C_BnD(\log n)^{12}$,
    put $a=W/K$. Then $Aa\sqrt B\le c$. If also $B\ge1$, then
    $Aa\le c/\sqrt B$.
\end{theorem}
\begin{proof}\leanok
    \uses{thm:al_moment_band_power,thm:al_moment_sqrt_budget}
    Choose $C_K$ before the scale. Theorems~\ref{thm:al_moment_band_power}
    and~\ref{thm:al_moment_sqrt_budget} give
    \begin{align}
        Aa\sqrt B
         & \le AC_K\sqrt{2C_B}\,n^{-\delta}(\log n)^6.
        \label{eq:al_moment_smallness}
    \end{align}
    Since $\delta>0$, the fixed multiple of $(\log n)^6$ is
    $o(n^\delta)$. Choose one threshold making the right side of
    (\ref{eq:al_moment_smallness}) at most $c$. This choice is independent
    of $W$ and $B$. For $B\ge1$, division by $\sqrt B>0$ proves the
    interval bound.
\end{proof}

\begin{theorem}[Every fixed multiple at the final collar scales]
    \label{thm:al_moment_final_scales}
    \lean{TNLean.PEPS.AreaLaw.Scan.final_scanner_moment_margin,
        TNLean.PEPS.AreaLaw.Scan.exists_final_a_le_moment_radius}
    \leanok
    Use the final parameters of
    \cite[Section~9, lines~839--865]{OpenAI2026AreaLaw}:
    \begin{align}
        \ell_{\rm s} & =1/100000,                  \\
        \mu          & =1/4,                       \\
        \kappa       & =1/1000000,                 \\
        \omega       & =1/20000,                   \\
        \nu          & =1/1000,\qquad e=2/1000000.
    \end{align}
    Their exact margin is
    \begin{align}
        1-\ell_{\rm s}-\mu-\omega-(1+\kappa)/2
         & =499879/2000000>0.
    \end{align}
    Fix $C_B\ge0$, $c>0$ and $t\in\mathbb R$. There is $N\ge2$ such
    that, for every $n\ge N$, the literal rounded scales in
    (\ref{eq:al_moment_rounded_scales}) have $K\ge1$ and, with
    $a=n^\omega/K$, every $1\le B\le C_BnD(\log n)^{12}$ satisfies
    $|ta|\le c/\sqrt B$.
\end{theorem}
\begin{proof}\leanok
    \uses{thm:al_moment_weight_smallness}
    Rational arithmetic gives the displayed positive margin and all
    stated exponent inequalities. Apply
    Theorem~\ref{thm:al_moment_weight_smallness} with $A=|t|$ and
    $W=n^\omega$. Since $a\ge0$, one has $|ta|=|t|a$.
\end{proof}
